\documentclass{article}
\usepackage[utf8]{inputenc}
% Asegura el uso correcto de caracteres especiales en español
\usepackage[spanish]{babel}
% Paquetes para notación matemática avanzada
\usepackage{amsmath}
\usepackage{amssymb}
% Define el formato de la página para márgenes razonables
\usepackage{geometry}
\geometry{a4paper, margin=1in}
% Para asegurar que el texto se justifique correctamente
\setlength{\parindent}{0pt} 

\begin{document}

\title{Examen de Matemáticas - Bachillerato}
\author{Operaciones Fundamentales y Álgebra}
\date{\today}
\maketitle

\section*{Instrucciones}
\begin{enumerate}
    \item Lea cuidadosamente cada problema antes de comenzar a resolverlo.
    \item Debe mostrar todo el procedimiento y los cálculos necesarios para obtener la puntuación completa.
    \item Responda las 5 preguntas prácticas.
\end{enumerate}

\vspace{0.5cm}

\section*{Preguntas Prácticas}

% Pregunta 1: Ecuaciones Lineales / Problema de Aplicación
\subsection*{Pregunta 1. Sistema de Ecuaciones Lineales}
Un cine vendió 150 boletos en una noche. Los boletos para adultos costaban $\$8$ y los boletos para niños costaban $\$5$. Si la recaudación total fue de $\$1020$, ¿cuántos boletos de cada tipo se vendieron? (Defina sus variables y plantee el sistema de ecuaciones).

\vspace{2.0cm} % Espacio para la respuesta

% Pregunta 2: Ecuación Cuadrática
\subsection*{Pregunta 2. Solución de Ecuación Cuadrática}
Resuelva la siguiente ecuación cuadrática. Muestre los pasos utilizando la fórmula general:
$$2x^2 - 3x - 2 = 0$$

\vspace{2.0cm}

% Pregunta 3: Simplificación de Polinomios
\subsection*{Pregunta 3. Simplificación y Expansión Algebraica}
Simplifique la siguiente expresión algebraica hasta su forma más reducida:
$$E(x) = (3x + 2)(x - 4) - (x^2 - 7x + 5)$$

\vspace{2.0cm}

% Pregunta 4: Evaluación de Funciones
\subsection*{Pregunta 4. Evaluación de una Función Racional y Radical}
Dada la función $f(x) = \frac{x^2 + 1}{\sqrt{2x + 15}}$, calcule el valor de $f(5)$. Exprese su respuesta final en forma de fracción simplificada.

\vspace{2.0cm}

% Pregunta 5: Teorema de Pitágoras y Razones Trigonométricas Básicas
\subsection*{Pregunta 5. Aplicación Trigonométrica en Triángulo Rectángulo}
Un constructor necesita calcular la longitud de una rampa que formará un ángulo de elevación $\alpha$ con el suelo. Si la rampa debe alcanzar una altura vertical de $3$ metros y la base horizontal que cubre es de $4$ metros, determine:
\begin{enumerate}
    \item La longitud (hipotenusa) de la rampa.
    \item El valor del $\cos(\alpha)$ (coseno del ángulo de elevación).
\end{enumerate}

\vspace{2.0cm}

\end{document}